% GENERATED FILE. Edit canonical lessons, then run npm run generate:books.
% canonical-source-hash: fnv1a64:4f5d39747cec9820
% canonical-lessons: GE-C95-heute, GE-C95-manchmal, GE-C95-dort, GE-C95-weit, GE-C95-rechts

\chapter{Today, Sometimes, There, Far, On the right}
\label{ch:ge-today-sometimes-there-far-on-the-right}
\hlchaptermodality{95}

\begin{chapteropening}
\textbf{By the end of this chapter:} \emph{I can say today, sometimes, there, far and on the right.}

Complete the last lesson of chapter 95: I can say today, sometimes, there, far and on the right.
\end{chapteropening}

\section[heute]{heute — today}

\label{lesson:GE-C95-heute}



\begin{quote}
Before the new one: say the German for always, then the German for immediately.
\end{quote}

\subsection*{You'll want to know: heute}
\textbf{heute} — \textquotedblleft{}today\textquotedblright{}.

\begin{cousinweb}
\textbf{heute} $\leftarrow$ Middle High German \emph{hiute} $\leftarrow$ Old High German \emph{hiutu}, \textquotedblleft{}on this day\textquotedblright{}, which ends in the old form of \emph{Tag} (EtymDB 2.1).
\end{cousinweb}

\begin{grammarlens}[title={What you've built}]
One more word for this chapter.
\end{grammarlens}

\subsection*{Your turn}
\begin{itemize}
\raggedright
  \item \emph{Say it:} \emph{heute}
  \item \emph{Say it:} \emph{heute}, once more
  \item \emph{Say it:} say \emph{heute}
  \item \emph{Recall:} say the German for always, then the German for immediately, then say \emph{heute} again
\end{itemize}

\begin{tcolorbox}[breakable,colback=teal!4,colframe=teal!35!black,title={Before you move on}]
What does heute mean? (\textquotedblleft{}Today\textquotedblright{}.) Say it once more.
\end{tcolorbox}
\section[manchmal]{manchmal — sometimes}

\label{lesson:GE-C95-manchmal}



\begin{quote}
Before the new one: say the German for immediately, then the German for today.
\end{quote}

\subsection*{You'll want to know: manchmal}
\textbf{manchmal} — \textquotedblleft{}sometimes\textquotedblright{}.

\begin{cousinweb}
\textbf{manchmal} is \textquotedblleft{}many a\textquotedblright{} + the word for \textquotedblleft{}a time\textquotedblright{}, Old High German \emph{māl} $\leftarrow$ Proto-Germanic \emph{*mēlą} (EtymDB 2.1): many a time.
\end{cousinweb}

\begin{grammarlens}[title={What you've built}]
One more word for this chapter.
\end{grammarlens}

\subsection*{Your turn}
\begin{itemize}
\raggedright
  \item \emph{Say it:} \emph{manchmal}
  \item \emph{Say it:} \emph{manchmal}, once more
  \item \emph{Say it:} say \emph{manchmal}
  \item \emph{Recall:} say the German for immediately, then the German for today, then say \emph{manchmal} again
\end{itemize}

\begin{tcolorbox}[breakable,colback=teal!4,colframe=teal!35!black,title={Before you move on}]
What does manchmal mean? (\textquotedblleft{}Sometimes\textquotedblright{}.) Say it once more.
\end{tcolorbox}
\section[dort]{dort — there}

\label{lesson:GE-C95-dort}



\begin{quote}
Before the new one: say the German for today, then the German for sometimes.
\end{quote}

\subsection*{You'll want to know: dort}
\textbf{dort} — \textquotedblleft{}there\textquotedblright{}.

\begin{cousinweb}
\textbf{dort} $\leftarrow$ Old High German \emph{dorot} (EtymDB 2.1) — inherited, built on the same pointing stem as English \emph{there}.
\end{cousinweb}

\begin{grammarlens}[title={What you've built}]
One more word for this chapter.
\end{grammarlens}

\subsection*{Your turn}
\begin{itemize}
\raggedright
  \item \emph{Say it:} \emph{dort}
  \item \emph{Say it:} \emph{dort}, once more
  \item \emph{Say it:} say \emph{dort}
  \item \emph{Recall:} say the German for today, then the German for sometimes, then say \emph{dort} again
\end{itemize}

\begin{tcolorbox}[breakable,colback=teal!4,colframe=teal!35!black,title={Before you move on}]
What does dort mean? (\textquotedblleft{}There\textquotedblright{}.) Say it once more.
\end{tcolorbox}
\section[weit]{weit — far}

\label{lesson:GE-C95-weit}



\begin{quote}
Before the new one: say the German for sometimes, then the German for there.
\end{quote}

\subsection*{You'll want to know: weit}
\textbf{weit} — \textquotedblleft{}far\textquotedblright{}.

\begin{cousinweb}
\textbf{weit} $\leftarrow$ Old High German \emph{wīt} $\leftarrow$ Proto-Germanic \emph{*wīdaz}, \textquotedblleft{}wide, far\textquotedblright{} — inherited, the same word as English \emph{wide}. (EtymDB 2.1)
\end{cousinweb}

\begin{grammarlens}[title={What you've built}]
One more word for this chapter.
\end{grammarlens}

\subsection*{Your turn}
\begin{itemize}
\raggedright
  \item \emph{Say it:} \emph{weit}
  \item \emph{Say it:} \emph{weit}, once more
  \item \emph{Say it:} say \emph{weit}
  \item \emph{Recall:} say the German for sometimes, then the German for there, then say \emph{weit} again
\end{itemize}

\begin{tcolorbox}[breakable,colback=teal!4,colframe=teal!35!black,title={Before you move on}]
What does weit mean? (\textquotedblleft{}Far\textquotedblright{}.) Say it once more.
\end{tcolorbox}
\section[rechts]{rechts — on the right}

\label{lesson:GE-C95-rechts}



\begin{quote}
Before the new one: say the German for there, then the German for far.
\end{quote}

\subsection*{You'll want to know: rechts}
\textbf{rechts} — \textquotedblleft{}on the right\textquotedblright{}.

\begin{cousinweb}
\textbf{rechts} is the word for \textquotedblleft{}right\textquotedblright{} + the adverb ending \emph{-s}. That word $\leftarrow$ Old High German \emph{reht} $\leftarrow$ Proto-Germanic \emph{*rehtaz}, the same word as English \emph{right}. (EtymDB 2.1)
\end{cousinweb}

\begin{grammarlens}[title={What you've built}]
One more word for this chapter.
\end{grammarlens}

\subsection*{Your turn}
\begin{itemize}
\raggedright
  \item \emph{Say it:} \emph{rechts}
  \item \emph{Say it:} \emph{rechts}, once more
  \item \emph{Say it:} say \emph{rechts}
  \item \emph{Recall:} say the German for there, then the German for far, then say \emph{rechts} again
\end{itemize}

\begin{tcolorbox}[breakable,colback=teal!4,colframe=teal!35!black,title={Before you move on}]
What does rechts mean? (\textquotedblleft{}On the right\textquotedblright{}.) Say it once more.
\end{tcolorbox}
